% Created 2026-05-06 Wed 17:37
% Intended LaTeX compiler: pdflatex
\documentclass[11pt]{article}
\usepackage[utf8]{inputenc}
\usepackage[T1]{fontenc}
\usepackage{graphicx}
\usepackage{longtable}
\usepackage{wrapfig}
\usepackage{rotating}
\usepackage[normalem]{ulem}
\usepackage{amsmath}
\usepackage{amssymb}
\usepackage{capt-of}
\usepackage{hyperref}
\author{fcalado}
\date{\today}
\title{}
\hypersetup{
 pdfauthor={fcalado},
 pdftitle={},
 pdfkeywords={},
 pdfsubject={},
 pdfcreator={Emacs 29.3 (Org mode 9.6.15)}, 
 pdflang={English}}
\begin{document}

\tableofcontents

\section{{\bfseries\sffamily TODO} Butler, Judith. \emph{Bodies That Matter}. 1993.}
\label{sec:org59e4da4}
\subsection{overview:}
\label{sec:org42119b9}
The book explores sex and gender "construction as a constitutive
constraint," (x) finds that it works by producing, as a side effect,
"a domain of unthinkable, abject, unlivable bodies" (x), and sets
about configuring this side effect as the raw material for bodily
resignifcation. For Butler, resignification of these citations is the
way out of the significatory circle.
\begin{itemize}
\item Bulter finds that \textbf{the body only materializes through language}. The
reality of eating, sleeping, fucking, etc, is "irrefutable." "But,"
says Butler, "their irrefutability in no way implies what it might
mean to affirm them and through what discursive means" (x).
\end{itemize}

She deploys examples of:
\begin{itemize}
\item Luce Irigaray's writing "citing" Plato;
\item kinship practices in film \emph{Paris Is Burning};
\item drag performance citing gender norms;
\item most famously, by the reclamation of the word "queer" citing its own
repudiation as an identity label.
\begin{itemize}
\item emphasizes the radical possibility of queer, how it reiterates
denigration to mean something else, a community with a future,
with value, "resignifying the abjection of homosexuality into
defiance and legitimacy" (xxviii).
\end{itemize}
\item reads aristotle, freud, lacan, etc, to tease out the
"indistinguishability of psychic and bodily formation" (xxix).
\end{itemize}

\subsection{key terms/concepts}
\label{sec:orgbe98338}
Resignifcation Process:
\begin{itemize}
\item The point isn't to demolish existing significatory systems, but to
rework them in a process of resignifcation.
\item Two steps:
\begin{enumerate}
\item Repetition: What we think of as bodily materiality is a
sedimentation that emerges from repetition, a kind of performance
whereby meaning is signified and re-signified endlessly, never
complete---citing whatever social power or understanding about
sex.
\begin{itemize}
\item Sex/gender are \textbf{unstable}, needing repetition:
\begin{itemize}
\item Sex/gender is never stable, that's why it needs repetition.
\end{itemize}
\item \textbf{Subjectivity} emerges in power structure:
\begin{itemize}
\item Subjects are always interpellated by the discourse prior to
citing it. The subject only comes into intelligibility
through the matrix of gender.
\end{itemize}
\end{itemize}
\item Re-formulating repudiation, as a troubling return:
\begin{itemize}
\item identification produces a repudiation, domain of abjection,
which always threatens the boundaries of the identification:
\begin{itemize}
\item "The forming of a subject requires an identification with the
normative phantasm of “sex,” and this identification takes
place through a repudiation which produces a \textbf{domain of
abjection}, a repudiation without which the subject cannot
emerge" (xiii).
\item "this disavowed abjection will \textbf{threaten to expose} the
self-grounding presumptions of the sexed subject, grounded as
that subject is in a repudiation whose consequences it cannot
fully control. The task will be to consider this threat and
disruption not as a permanent contestation of social norms
condemned to the pathos of perpetual failure, but rather as a
critical resource in the struggle to rearticulate the very
terms of symbolic legitimacy and intelligibility" (xiii).
\end{itemize}
\item The only freedom that is possible resides within this power of
discourse, resignifying it, perhaps through parody or
impersonation.
\begin{itemize}
\item Agency:
\item “The compulsion to repeat an injury is not necessarily the
compulsion to repeat the injury in the same way or to stay
fully within the traumatic orbit of that injury. The force of
repetition in language may be the paradoxical condition by
which a certain agency---not linked to a fiction of the ego
as master of circumstance---is derived from the impossibility
of choice…. Paris is Burning might be understood as
repetitions of hegemonic forms of power that fail to repeat
loyally and, in that failure, open possibilities for
resignifying the terms of violation against their violating
aims” (383).
\item "My purpose here is to understand how what has been
foreclosed or banished from the proper domain of "sex" --
where that domain is secured through a heterosexualizing
imperative -- might at once be produced as a \textbf{troubling
return}, not only as an \emph{imaginary} contestation that effects
a failure in the workings of the inevitable law, but as an
enabling disruption, the occasion for a radical
rearticulation of the symbolic horizon in which bodies come
to matter at all" (XXX).
\end{itemize}
\item Examples of resignification that draws from repudiation:
\begin{itemize}
\item We see this in the word “queer” which has been
re-appropriated---something thatsignified abjectness now
means defiance. We can also use repetitionto re-signify
identification, to the point where it loses its power.
\item In Paris is Burning, not only are the male drag
performersexposing the superficiality of gender, but also
performing care in away that is feminine, “mothering”
“housing” “rearing” each other.
\end{itemize}
\end{itemize}
\end{enumerate}
\end{itemize}

Language:
\begin{itemize}
\item language is productive, can resignify meaning. It is the option
available to those who are trapped within the signification system.
\begin{itemize}
\item "The body posited as prior to the sign, is always posited or
signified as prior. This signification produces as an effect of
its own procedure the very body that it nevertheless and
simultaneously claims to discover as that which precedes its own
action. If the body signifi ed as prior to signifi cation is an
effect of signification, then the mimetic or representational
status of language, which claims that signs follow bodies as their
necessary mirrors, is not mimetic at all. On the contrary, it is
productive, constitutive, one might even argue performative,
inasmuch as this signifying act delimits and contours the body
that it then claims to fi nd prior to any and all signification."
(6).
\end{itemize}
\end{itemize}

On proliferation:


\subsection{close readings}
\label{sec:org31bfc35}
On \emph{Paris is Burning}:
\begin{itemize}
\item "This “being a man” and this “being a woman” are internally
unstable affairs. They are always beset by ambivalence precisely
because there is a cost in every identification, the loss of some
other set of identifications, the forcible approximation of a norm
one never chooses, a norm that chooses us, but which we occupy,
reverse, resignify to the extent that the norm fails to determine
us completely." (86).
\item when resignification doesn't always work, Venus Xtravaganza:
\begin{itemize}
\item "The citing of the dominant norm does not, in this instance,
displace that norm; rather, it becomes the means by which that
dominant norm is most painfully reiterated as the very desire
and the performance of those it subjects\ldots{} As much as she
crosses gender, sexuality, and race performatively, the hegemony
that reinscribes the privileges of normative femininity and
whiteness wields the fi nal power to re naturalize Venus’s body
and cross out that prior crossing, an erasure that is her death"
(91).
\end{itemize}
\item Kinship norms are resignified:
\begin{itemize}
\item "it is in the elaboration of kinship forged through a resignifi
cation the very terms which effect our exclusion and abjection
that such a resignifi cation creates the discursive and social
space for community, that we see an appropriation of the terms of
domination that turns them toward a more enabling future." (95).
\end{itemize}
\end{itemize}

Drag and heterosexual melancholia:
\begin{itemize}
\item "that drag exposes or allegorizes the mundane psychic and
performative practices by which heterosexualized genders form
themselves through the renunciation of the possibility of
homosexuality, a foreclosure that produces a fi eld of heterosexual
objects at the same time that it produces a domain of those whom it
would be impossible to love. Drag thus allegorizes heterosexual
melancholy , the melancholy by which a masculine gender is formed
from the refusal to grieve the masculine as a possibility of love; a
feminine gender is formed (taken on, assumed) through the
incorporative fantasy by which the feminine is excluded as a
possible object of love, an exclusion never grieved, but “preserved”
through the heightening of feminine identification itself. In this
sense, the “truest” lesbian melancholic is the strictly straight
woman, and the “truest” gay male melancholic is the strictly
straight man" (180).
\item "The critical promise of drag does not have to do with the
proliferation of genders, as if a sheer increase in numbers would do
the job, but rather with the exposure or the failure of heterosexual
regimes ever fully to legislate or contain their own ideals" (181).
\end{itemize}

On Irigaray:
\begin{itemize}
\item "For how can one read a text for what does \emph{not} appear within its
own terms, but which nevertheless constitutes the illegible
conditions of its own legibility? (11)
\item One cannot interpret the philosophical relation to the feminine
through the figures that philosophy provides, but, rather, she
argues, through siting [sic] the feminine as the unspeakable
condition of figuration, as that which, in fact, can never be
figured within the terms of philosophy proper, but whose exclusion
from that propriety is its enabling condition" (12).
\begin{itemize}
\item "This is a taking of his [Plato's] place, not to assume it, but to
show that it is \emph{occupiable}\ldots{} a kind of overreading which mimes
and exposes the speculative excesses in Plato (emphasis original,
11).
\end{itemize}
\item "This explains in part the radical citational practice of Irigaray,
the catachrestic usurpation of the “proper” for fully improper
purposes. For she mimes philosophy—-as well as psychoanalysis—-and,
in the mime, takes on a language that effectively cannot belong to
her, only to call into question the exclusionary rules of
proprietariness that govern the use of that discourse. This
contestation of propriety and property is precisely the option open
to the feminine when it has been constituted as an excluded
impropriety, as the improper, the propertyless" (12).
\item "Irigaray’s response to this exclusion of the feminine from the
economy of representation\ldots{} I’ll show you what this unintelligible
receptacle can do to your system; I will not be a poor copy in your
system, but I will resemble you nevertheless by miming the textual
passages through which you construct your system and showing that
\textbf{what cannot enter it is already inside it (as its necessary
outside)}, and I will mime and repeat the gestures of your operation
until this emergence of the outside within the system calls into
question its systematic closure and its pretension to be
self-grounding\ldots{}.In this sense, she performs a repetition and
displacement of the phallic economy. This is citation, not as
enslavement or simple reiteration of the original, but as an
insubordination that appears to take place within the very terms of
the original, and which calls into question the power of origination
that Plato appears to claim for himself. Her miming has the effect
of repeating the origin only to displace that origin as an origin"
(18).
\item On "matter" association to femininity:
\begin{itemize}
\item "Irigaray’s task\ldots{} is to show that those binary oppositions are
formulated through the exclusion of a field of disruptive
possibilities. Her speculative thesis is that those binaries, even
in their reconciled mode, are part of a phallogocentric economy
that produces the “feminine” as its constitutive outside.
Irigaray’s intervention in the history of the form/matter
distinction underscores “matter” as the site at which the feminine
is excluded from philosophical binaries" (10).
\item "The economy that claims to include the feminine as the
subordinate term in a binary opposition of masculine/feminine
excludes the feminine, produces the feminine as that which must be
excluded for that economy to operate" (10).
\item "every explicit distinction takes place in an inscriptional space
that the distinction itself cannot accomodate. Matter as a \emph{site}
of inscription cannot be explicitly thematized. And this
inscriptional site or space is, for Irigaray, a \emph{materiality} that
is not the same category of "matter" whose articulation it
conditions and enables\ldots{} this second inarticulate "matter"
designates the constitutive outside of the Platonic economy; it is
what must be excluded for that economy to posture as internally
coherent" (12-13).
\end{itemize}

\item A fount for future subversions:
\begin{itemize}
\item "The task is to refigure this necessary “outside” as a future
horizon, one in which the violence of exclusion is perpetually in
the process of being overcome. But of equal importance is the
preservation of the outside, the site where discourse meets its
limits, where the opacity of what is not included in a given regime
of truth acts as a disruptive site of linguistic impropriety and
unrepresentability, illuminating the violent and contingent
boundaries of that normative regime precisely through the inability
of that regime to represent that which might pose a fundamental
threat to its continuity. In this sense, \textbf{radical and inclusive
representability is not precisely the goal}: to include, to speak
as, to bring in every marginal and excluded position within a given
discourse is to claim that a singular discourse meets its limits
nowhere, that it can and will domesticate all signs of difference.
If there is a violence necessary to the language of politics, then
the risk of that violation might well be followed by another in
which we begin, without ending, without mastering, to own—-and yet
never fully to own—-the exclusions by which we proceed." (25).
\item She (like Munoz, Love?) is looking toward the horizon, where
queerness can never be defined, but must remain the outside that
delimits the known. The goal is not to be inclusive, it's to
preserve this space that cannot be included without being coopted.
\end{itemize}
\end{itemize}

\subsection{critiques}
\label{sec:org8ceab4e}
\subsubsection{Jay Prosser}
\label{sec:org41276bf}
\begin{itemize}
\item Bulter finds that the body only materializes through language/norms.
The reality of eating, sleeping, fucking, etc, is "irrefutable."
"But," says Butler, "their irrefutability in no way implies what it
might mean to affirm them and through what discursive means" (x).
\end{itemize}
\subsubsection{Martha Nussbaum}
\label{sec:org6af2bf2}
\begin{itemize}
\item Thinks Butler's politics are ineffective, weak, doesn't serve women.
\begin{itemize}
\item "A new and disquieting trend" in feminism\ldots{} "virtually complete
turning from the material side of life, toward a type of verbal
and symbolic politics that makes only the flimsiest of connections
with the real situation of women."
\end{itemize}
\item Nussbaum claims that Butler's position encourages feminists to
"adopt a stance that looks like quietism and retreat".
\item Critique of writing style, "teasing, exasperating" style that
"bullies the reader". Not philosophy, but "sophistry and rhetoric".
\begin{itemize}
\item Butler's style of posing questions - "mystification as well as
hierarchy are the tools of her practice." "Butler never quite
tells the reader whether she approves of the view described."
\end{itemize}
\item Biggest critique: that parody and subversion are weak tools:
\begin{itemize}
\item "What exactly does she council when she councils subversion?"
\item Slavery wasn't overcome through parody.
\item Parody won't work for women who are "hungry, illiterate,
disenfranchised, beaten, raped\ldots{} Such women prefer schools,
votes, and the integrity of their bodies."
\end{itemize}
\end{itemize}
\end{document}
